\documentclass[11pt,a4paper]{article}
\usepackage[T1]{fontenc}
\usepackage[utf8]{inputenc}
\usepackage{amsmath,amssymb,amsthm}
\usepackage[margin=1in]{geometry}
\usepackage[colorlinks=true,linkcolor=blue,urlcolor=blue]{hyperref}
\theoremstyle{plain}
\newtheorem{theorem}{Theorem}
\newtheorem{lemma}[theorem]{Lemma}
\newtheorem{proposition}[theorem]{Proposition}
\newtheorem{corollary}[theorem]{Corollary}
\theoremstyle{definition}
\newtheorem{remark}[theorem]{Remark}

\title{The cross-base identity for OEIS A126416, proved}
\author{Adrian Perez Fontelles\\ \small Independent researcher}
\date{13 September 2026}
\begin{document}
\maketitle

\begin{abstract}
OEIS A126416 carries an empirical identity, in two clauses, relating its counts to those of the
entry one base smaller, contributed by R. H. Hardin. Both clauses are true and both are proved
here, with no computation in the argument. The difference between consecutive bases counts the
admissible strings that USE the new largest digit; those are the translates of the
translation-classes of admissible strings over $\mathbf{Z}$ whose range fits below the top of
the alphabet; the classes are the step vectors, of which there are $(2d+1)^{n-1}$; and at the
one base where the fit fails by exactly one, the classes lost are the two monotone ones, which
is the entry's own $-2$.
\end{abstract}

\noindent\small 2020 Mathematics Subject Classification. 05A15, 05A05, 05C38.\normalsize

\section{The conjecture}

OEIS A126416 is ``Number of base 29 n-digit numbers with adjacent digits differing by two or less.''. It has offset $0$ and begins
\[
1, 29, 139, 675, 3293, 16111, 78973, 387649,\ \dots
\]
The entry states, as an empirical claim and never marked settled:
\begin{quote}
[Empirical] a(base,n)=a(base-1,n)+5\textasciicircum{}(n-1) for base>=2n-1; a(base,n)=a(base-1,n)+5\textasciicircum{}(n-1)-2 when base=2n-2
\end{quote}
As of the ``Last modified'' line on the live entry (Aug 12 2026 19:12:32, revision 9)
this is still recorded as empirical, and nothing on the entry records it as proved.

For this entry the base is $29$ and the permitted difference is $d=2$, so the first clause
claims $a_{29}(n)-a_{28}(n)=5^{\,n-1}$ whenever $2(n-1)+1\le29$, that is for
$1\le n\le15$, and the second claims $a_{29}(n)-a_{28}(n)=5^{\,n-1}-2$ at the single
$n$ with $2(n-1)=29$.

\section{What the entry counts}

A base-$b$ $n$-digit string in this family's sense is a tuple
$(d_1,\dots,d_n)\in\{0,1,\dots,b-1\}^{n}$ with $|d_i-d_{i+1}|\le d$ for $1\le i<n$; there
is no wrap-around. Leading zeros are counted, which is what the entry's $a(1)=b$ says, and
$a(0)=1$ is the empty string. Write $L_d(b,n)$ for their number. Evaluated exactly in integer
arithmetic, $L_{2}$ at base $29$ reproduces all $21$ terms A126416 publishes.

\section{The identity}

\begin{theorem}
Fix $d\ge1$ and $n\ge1$, and put $K=2d+1$.
\begin{enumerate}
\item[(i)] If $b\ge d(n-1)+1$ then $L_d(b,n)-L_d(b-1,n)=K^{\,n-1}$.
\item[(ii)] If $b=d(n-1)$ then $L_d(b,n)-L_d(b-1,n)=K^{\,n-1}-2$.
\end{enumerate}
\end{theorem}

\begin{proof}
\emph{Step 1.} The admissible strings over $\{0,\dots,b-1\}$ that do not use the value $b-1$
are exactly the admissible strings over $\{0,\dots,b-2\}$, so $L_d(b,n)-L_d(b-1,n)$ is the
number of admissible strings over $\{0,\dots,b-1\}$ in which $b-1$ occurs.

\emph{Step 2.} Consider the admissible strings over $\mathbf{Z}$: tuples
$(c_1,\dots,c_n)\in\mathbf{Z}^{n}$ with $|c_i-c_{i+1}|\le d$. Translation by an integer acts
freely on them, each orbit contains exactly one representative with $c_1=0$, and such a
representative is determined by its step vector $(s_1,\dots,s_{n-1})$, $s_i=c_{i+1}-c_i$,
which may be any element of $\{-d,\dots,d\}^{n-1}$. So there are exactly $K^{\,n-1}$ orbits.
Each orbit also contains exactly one representative with $\max_i c_i=0$, and all its
representatives share one RANGE $\rho=\max_ic_i-\min_ic_i$, which satisfies
$\rho\le d(n-1)$ because each of the $n-1$ steps moves by at most $d$.

\emph{Step 3.} A string counted in Step 1 has maximum $b-1$ and entries at least $0$, so it is
the representative with maximum $b-1$ of an orbit of range at most $b-1$; conversely every orbit
of range at most $b-1$ gives exactly one such string. Hence
\[
L_d(b,n)-L_d(b-1,n)\;=\;\#\{\text{orbits of range}\le b-1\} .
\]

\emph{Step 4.} If $b\ge d(n-1)+1$ then $b-1\ge d(n-1)$ and by Step 2 every orbit qualifies,
giving (i). If $b=d(n-1)$ then exactly the orbits of range $d(n-1)$ are excluded. An orbit has
range $d(n-1)$ only if some $c_i$ attains the maximum, some $c_j$ the minimum, and the walk
between them covers $d(n-1)$; that walk has at most $n-1$ steps, each of size at most $d$, so it
must use all $n-1$ of them in the same direction at full size. There are exactly two such step
vectors, $(d,\dots,d)$ and $(-d,\dots,-d)$, so two orbits are lost and (ii) follows.
\end{proof}

\section{Verification}

Three checks, each independent of the proof above.

First, the model reproduces all $21$ terms A126416 publishes, exactly, in integer
arithmetic; the reading of the entry's English is what is being tested there.

Second, the first clause was evaluated directly for every $n$ in the entry's claimed range,
$1\le n\le15$, both sides computed from the definition and not from the theorem:

\begin{center}
\begin{tabular}{rrrr}
$n$ & $L_{2}(29,n)$ & $L_{2}(28,n)$ & $5^{\,n-1}$ \\ \hline
1 & 29 & 28 & 1 \\
2 & 139 & 134 & 5 \\
3 & 675 & 650 & 25 \\
4 & 3293 & 3168 & 125 \\
5 & 16111 & 15486 & 625 \\
6 & 78973 & 75848 & 3125 \\
7 & 387649 & 372024 & 15625 \\
8 & 1904855 & 1826730 & 78125 \\
\end{tabular}
\end{center}

\noindent and the difference of the second and third columns is the fourth at every row shown
and at every row up to $n=15$.

Third, the second clause was evaluated at its own index.

\section*{References}
\begin{enumerate}
\item The OEIS Foundation Inc., \emph{The On-Line Encyclopedia of Integer Sequences},
\url{https://oeis.org/A126416}.
\end{enumerate}

\end{document}
